\documentclass[a4paper]{article}
\usepackage[pdftex]{graphicx}
\usepackage[utf8]{inputenc}
\usepackage{enumerate}
\usepackage{mdwlist}
\usepackage{amssymb}
\usepackage{icomma}
\usepackage{siunitx}
\sisetup{locale=DE} 
\usepackage{href-ul}
\hypersetup{
	colorlinks=true,
	linkcolor=blue,
	urlcolor=blue}
\usepackage{geometry}
\geometry{a4paper, top=15mm, left=15mm, right=15mm, bottom=15mm,
	headsep=10mm, footskip=12mm}

\begin{document}
	{\bf Freier Fall und senkrechter Wurf}

\begin{enumerate}[1.]
{\bf \item Der Legende nach lag Sir Isaac Newton (1643–1727) } unter einem Apfelbaum, als er einen Apfel beim Fallen beobachtete und daraus auf 
sein Kraftgesetz schloss. Wie viel Reaktionszeit hätte er gehabt, den Apfel aufzufangen, wenn die Fallhöhe $h=\SI{2,4}{\meter}$ gewesen wäre?
Wie schnell wäre der Apfel dann geworden?

{\bf \item Eine Paketdrohne kommt geflogen und liefert} ein Päckchen für Frau Eilig aus, die schon vor ihrem Haus steht.  Wie hoch muss die Flughöhe der Drohne minimal sein, aus welcher sie das Päckchen abwirft, wenn Frau Eilig $\SI{1,70}{\meter}$ groß ist und $\SI{0,5}{\second}$ Reaktionszeit   benötigt, um das Päckchen zu fangen? Wie viel kinetische Energie hat das $\SI{1,5}{\kilogram}$ schwere Päckchen beim Fangen?  

{\bf \item Das Wasser der Fontäne im Künettegraben} strömt mit einer Geschwindigkeit von $\SI{9,0}{\meter\over\second}$ aus. Wie hoch ist die Fontäne?
Wie lange ist ein Wassertropfen des Strahls in der Luft, bevor er wieder unten aufkommt?

{\bf \item TEXUS -- Technische Experimente Unter Schwerelosigkeit} Hierzu werden Raketen in Kiruna/Nordschweden senkrecht nach oben abgeschossen. Die Nutzlast ist beim ballistischen Flug dabei $\SI{6,0}{\minute}$ nahezu unter Schwere\-losigkeit.
\begin{enumerate}[(a)]
	\item Welche Flughöhe wird dabei theoretisch benötigt? 
	\item Tatsächlich ist die erreichte Flughöhe $\SI{250}{\kilo\meter}$. Nenne hierfür zwei Gründe! 
\end{enumerate} 

{\bf \item Leider ging der Elfmeter an die Latte und der Ball prallt senkrecht nach oben ab.} Die Spidercam befindet sich direkt $\SI{3,0}{\meter}$ darüber und filmt das ganze. 
\begin{enumerate}[(a)]
	\item Wird die Kamera getroffen, wenn der Ball vor dem Abprall $v=\SI{11}{\meter\over\second}$ schnell war und durch den Abprall 
	50\% seiner kinetischen Energie verliert?  (Probiere, ohne die Masse zu rechnen -- es ist möglich!)
\suspend{enumerate}
Bei einem zweitem Elfmeter wird der Ball wieder gleich stark und wieder an die Latte geschossen, nur prallt er diesmal nach unten ab .
\resume{enumerate}
	\item Welche kinetische Energie hat er beim Aufprall auf den am Boden liegenden Torwart, wenn die Latte $\SI{2,4}{\meter}$ über dem Boden ist und die Masse des Balls $m=\SI{0,42}{\kilogram}$ beträgt ? 
	\item Welcher Geschwindigkeit entspricht das? Rechne noch einmal ohne Masse!
\end{enumerate} 


{\bf \item Lunerald steht auf dem Drei-Meter-Brett.} Sein Körperschwerpunkt befindet sich zunächst auf einer Höhe von $h=\SI{4,0}{\meter}$ über dem Wasser und seine Masse beträgt $m=\SI{75}{\kilogram}$. Beim Absprung dehnt sich das Brett mit der Federhärte $D=\SI{7,5}{\kilo\newton\over\meter}$
um $s=\SI{0,40}{\meter}$
\begin{enumerate}[a)]
	\item Wie hoch ist die Spannenergie beim Absprung ?
	\item Welche maximale Höhe über dem Wasser erreicht Lunerald ?
	\item Wie viel Zeit hat er für einen Salto ? (=Flugdauer)
\end{enumerate} 

{\bf \item Estela ist bereit zum Bungee-Sprung.} Sie steht auf einem  $h_k=\SI{80}{\meter}$ hohen Kran, das elastische Seil ist ungedehnt $l_s=\SI{30}{\meter}$ lang. Zur Sicherheit soll der Fall $\SI{15}{\meter}$ über dem Boden gestoppt werden.
\begin{enumerate}[a)]
	\item Wie lange ist Dauer des freien Falls und welche Maximalgeschwindigkeit erreicht sie dabei?
    \item Welche Höhe $h_r$ über dem Boden erreicht sie beim ,,Rebound", also beim Wiederaufstieg durch das Gummiseil, wenn durch Luftwiderstand und Verformungsarbeit 40\% der Energie verloren geht?
    \item Welchen Betrag muss die Federhärte des Seils mindestens sein, wenn Estelas Masse $m=\SI{60}{\kilogram}$ beträgt?
\end{enumerate} 
\fbox{
	\begin{minipage}{0.5\textwidth}
		Zur Lösung bitte \href{https://www.okuyakl.de/phys/p9frfL075/ll075.pdf}{hier klicken} oder den QR-Code scannen.\\
	Weitere Arbeitsblätter gibt es unter 
	
	\href{https://www.okuyakl.de}{www.okuyakl.de}
	\end{minipage}
	\hfill
	\begin{minipage}{0.4\textwidth}
		\includegraphics[width=1.5 cm]{../../viecher/zwe03}
		\includegraphics[width=3 cm]{qr075}
		\includegraphics[width=2 cm]{../../viecher/afanticon1}
		
\end{minipage}}

\end{enumerate} 

\end{document}

